\subsection{Fixed-estimator counterfactual mixture development}
The fixed-estimator CDM study has completed all nine source/seed pairs.
The initial seven-pair generation and its two originally unavailable MUCars
pairs retain separate immutable provenance; their complete evidence union
has an independent audit of 1,512 execution cells and 108 calibration fits.
Every stored payoff row, full-matrix LP optimum, rule allocation and mixture
expectation is checked. Table~\ref{tab:cdm_development} includes all three
scheduled sources and all three fitting seeds.

\begin{table*}[t]
\centering
\caption{CDM development at 10\% capacity with fixed estimators. Entries are
three-seed means of positive expected-log uplift. The final column is the
conditional shared-batch 95\% interval for risk minus CDM; positive values
favor CDM. This compares verification rules within each fixed estimator,
rather than a matched full-estimator comparison with feedback-trained SAV.}
\label{tab:cdm_development}
\begin{tabular}{llrrrrl}
\toprule
Source & Estimator & CDM & Risk & Benefit & Uniform & Risk $-$ CDM interval\\
\midrule
Synthetic & Mask-aware & 0.06224 & 0.05634 & 0.06152 & 0.07968 & $[-0.00779,-0.00385]$\\
 & Robust & 0.01192 & 0.01005 & 0.01017 & 0.01340 & $[-0.00291,-0.00110]$\\
\midrule
JUCars & Mask-aware & 0.01879 & 0.01854 & 0.01953 & 0.01886 & $[-0.00042,-0.00012]$\\
 & Robust & 0.01063 & 0.01055 & 0.01126 & 0.01009 & $[-0.00044,0.00029]$\\
\midrule
MUCars & Mask-aware & 0.04398 & 0.03102 & 0.04605 & 0.04430 & $[-0.01461,-0.01144]$\\
 & Robust & 0.00699 & 0.00626 & 0.00782 & 0.00658 & $[-0.00102,-0.00045]$\\
\bottomrule
\end{tabular}
\end{table*}

On the synthetic cohort, CDM reduces uplift relative to uniform mixing but
has higher uplift than the simple risk rule under both estimators. On JUCars,
it improves on benefit priority but does not establish a risk-rule advantage;
under the robust estimator it also loses to uniform mixing. All twelve pure
rules, the training-selected pure rule, both mixtures and all six capacities
are retained in the accompanying contrast data. These comparisons with uniform mixing do not establish an advantage
over the stronger risk rule.

The complete MUCars replication also favors risk priority over CDM under
both price predictors: CDM uplift is 41.79\% and 11.63\% higher for
mask-aware and robust prices. The corresponding risk-minus-CDM intervals
are below zero. With robust prices, CDM improves on benefit priority by
10.62\% and on the training-selected pure rule by 6.55\%, but has 6.32\%
more uplift than uniform mixing. At the preset 30\% capacity, robust-price
MUCars instead has 22.19\% less CDM uplift than risk priority, with paired
difference $0.00136$ and exploratory interval $[0.00107,0.00169]$.
This capacity-dependent result does not replace the primary comparison.
A finite calibration improvement does not imply a general held-out average
improvement over simple alternatives.

\begin{figure*}[t]
\centering
\includegraphics[width=\textwidth]{figures/cdm_capacity_frontiers_full.pdf}
\caption{Complete CDM source study at every preset capacity. Top:
mask-aware price predictor; bottom: robust/no-feedback price predictor.
Points average three fitting seeds and bars give their sample standard
deviations. The 10\% primary capacity is marked. All sources and capacities
remain visible; no favorable rule or source is selected after evaluation.
Conditional risk-minus-CDM primary intervals appear in
Table~\ref{tab:cdm_development}.}
\label{fig:cdm_capacity}
\end{figure*}

The finite LP guarantee concerns the stored calibration types and peer
profiles. It does not guarantee a reduction of mean held-out uplift, and its
worst-gain objective differs from that average. In the source experiment,
the honest-error constraint is invariant across rules and contributes no
observed discrimination. The mechanism's usefulness therefore remains an
empirical question despite exact finite optimization.
